\BourbakiExerciseGroup

\begin{enumerate}
\def\labelenumi{\arabic{enumi})}
\item
  Find all possible filters on a finite set.
\item
  If the intersection of all the sets of a filter F on an infinite set X is empty, show that F is finer than the filter of complements of finite subsets of X.
\item
  The intersection of two filters \(\mathfrak{F}\) and \(\mathfrak{G}\) on a set \(X\) is the set of all \(M \cup N\) , where \(M \in \mathfrak{F}\) and \(N \in \mathfrak{G}\) .
\item
  Let X be an infinite set. Show that the filter of complements of finite subsets of X is the intersection of the elementary filters associated with infinite sequences in X all of whose terms are distinct.
\item
  Let \(\mathfrak{F}\) and \(\mathfrak{G}\) be two filters on a set X. Show that if \(\mathfrak{F}\) and \(\mathfrak{G}\) have a least upper bound in the set of all filters on X, then this least upper bound is the set of all \(M \cap N\) , where \(M \in \mathfrak{F}\) and \(N \in \mathfrak{G}\) .
\item
  If to each filter on a set X we make correspond the associated topology (no. 5) on the set \(X' = X \cup \{\omega\}\) , then we have the following propositions:
\end{enumerate}

\begin{enumerate}
\def\labelenumi{\alph{enumi})}
\item
  If a topology \(\mathfrak{C}\) on \(\mathbf{X}'\) is finer than the topology associated with a filter \(\mathfrak{F}\) on \(\mathbf{X}\) , then \(\mathfrak{C}\) is either the discrete topology or the topology associated with a filter finer than \(\mathfrak{F}\) . Converse?
\item
  The greatest lower bound of the topologies associated with the filters of a set \(\Phi\) of filters on X is the topology associated with the intersection of the filters belonging to \(\Phi\) .
\item
  Let \(\mathfrak{G}\) be a set of subsets of X, and let \(\mathfrak{G}'\) be the set of subsets of X' of the form \(\mathbf{M} \cup \{\omega\}\) where \(\mathbf{M} \in \mathfrak{G}\) . Let \(\mathfrak{G}\) denote \(\mathfrak{G} \cup \mathfrak{P}(\mathbf{X})\) . If \(\mathfrak{G}\) generates a filter \(\mathfrak{F}\) , show that \(\mathfrak{G}\) is a subbase of the topology on X' associated with \(\mathfrak{F}\) . What topology does \(\mathfrak{G}\) generate if \(\mathfrak{G}\) is not a filter subbase?
\item
  \(\mathfrak{F}\) is an ultrafilter on \(\mathbf{X}\) if and only if the associated topology on \(\mathbf{X}'\) is \(\mathbf{X}\) -maximal ( \(\S 3\) , Exercise 11).
\end{enumerate}

\begin{enumerate}
\def\labelenumi{\arabic{enumi})}
\setcounter{enumi}{6}
\item
  In a topological space X, the intersection of the neighbourhood filters of all the points of a non-empty subset A of X is the neighbourhood filter of A in X.
\item
  \begin{enumerate}
  \def\labelenumii{\alph{enumii})}
  \tightlist
  \item
    Let \(\Phi\) be a set of topologies on a set \(X\) . Show that the neighbourhood filter of a point \(x \in X\) with respect to the intersection of the topologies of \(\Phi\) is coarser than the intersection of the neighbourhood filters of \(x\) with respect to the topologies belonging to \(\Phi\) .
  \end{enumerate}
\end{enumerate}

\begin{enumerate}
\def\labelenumi{\alph{enumi})}
\setcounter{enumi}{1}
\tightlist
\item
  Let \(\mathfrak{G}_1\) denote the topology \(\mathfrak{G}_0(\mathbf{Q}^2)\) (§ 2, Exercise 7) where the set \(\mathbf{Q}^2\) is linearly ordered by the lexicographic ordering (Set Theory,
\end{enumerate}

Chapter III, § 2, no. 6). Let \(\mathcal{G}_2\) denote the topology on \(\mathbf{Q}^2\) obtained by transporting \(\mathcal{G}_1\) by means of the symmetry \((\xi, \eta) \to (\eta, \xi)\) . Show that if \(\mathcal{G}\) is the intersection of the topologies \(\mathcal{G}_1, \mathcal{G}_2\) , then the neighbourhood filter of a point of \(\mathbf{Q}^2\) with respect to \(\mathcal{G}\) is strictly coarser than the intersection of the neighbourhood filters of the point with respect to \(\mathcal{G}_1\) and \(\mathcal{G}_2\) .

\begin{enumerate}
\def\labelenumi{\arabic{enumi})}
\setcounter{enumi}{8}
\tightlist
\item
  \begin{enumerate}
  \def\labelenumii{\alph{enumii})}
  \tightlist
  \item
    Show that every ultrafilter which is finer than the intersection of a finite number of filters is finer than at least one of them (use Proposition 5 of no. 4).
  \end{enumerate}
\end{enumerate}

\begin{enumerate}
\def\labelenumi{\alph{enumi})}
\setcounter{enumi}{1}
\tightlist
\item
  Give an example of an ultrafilter which is finer than the intersection of an infinite family of ultrafilters, but which is not equal to any of the ultrafilters of this family (consider the family of ultrafilters on an infinite set X each of which has a base consisting of a single point).
\end{enumerate}

\begin{enumerate}
\def\labelenumi{\arabic{enumi})}
\setcounter{enumi}{9}
\item
  Show that the intersection of all the sets of an ultrafilter contains at most one point; if one point, then the ultrafilter consists of all the sets containing this point (use Proposition 5 of no. 4).
\item
  Show that if a subset A of a set X does not belong to an ultrafilter U on X, then the trace of U on A is the set \(\mathfrak{P}(A)\) .
\item
  On an infinite set, show that an elementary filter associated with a sequence all of whose terms are distinct is not an ultrafilter.
\item
  Let \(f\) be a mapping of a set \(\mathbf{X}\) into a set \(\mathbf{X}'\) . Then \(f\) is injective if and only if, for every filter base \(\mathfrak{B}\) on \(\mathbf{X}\) , \((f\overline{f}^{1}(\mathfrak{B}))\) is a filter base equivalent to \(\mathfrak{B}\) .
\item
  Show that if \(f\) is a mapping of a set \(\mathbf{X}\) onto a set \(\mathbf{X}'\) , then the image \(f(\mathfrak{F})\) of every filter \(\mathfrak{F}\) on \(\mathbf{X}\) is a filter on \(\mathbf{X}'\) .
\end{enumerate}

¶ 15) Let \(n \to f(n)\) be a mapping of \(\mathbf{N}\) onto \(\mathbf{N}\) such that \(\overline{f}(m)\) is finite for each \(m \in \mathbf{N}\) . If \((x_n)\) is any sequence in a set \(\mathbf{X}\) , put \(y_n = x_{f(n)}\) . Show that the elementary filters associated with the sequences \((x_n)\) and \((y_n)\) are the same.

Deduce that if \((a_{n})\) and \((b_{n})\) are two sequences of elements of a set X such that the filter associated with \((b_{n})\) is finer than the filter associated with \((a_{n})\) , then the filter associated with \((b_{n})\) is the same as the filter associated with some subsequence of \((a_{n})\) .

¶ 16) Let \(\Phi\) be a countable linearly ordered set of elementary filters on a set X. Show that there is an elementary filter which is finer than all the filters of \(\Phi\) (show that the union of the filters of \(\Phi\) has a countable base).

\begin{enumerate}
\def\labelenumi{\arabic{enumi})}
\setcounter{enumi}{16}
\tightlist
\item
  Let X be a lattice (Set Theory, Chapter III, § 1, no. 13). A non-empty subset F of X is called a prefilter if it satisfies the following conditions: (i) whenever \(x \in \mathbf{F}\) and \(y \geqslant x\) , then \(y \in \mathbf{F}\) ; (ii) whenever \(x \in \mathbf{F}\) and \(y \in F\) , then \(\inf (x, y) \in F\) ; (iii) \(F \neq X\) . Thus the prefilters on the set of subsets of a set Y, ordered by inclusion, are precisely the filters on Y. A prefilter F is said to be prime if the relation \(\sup(x, y) \in F\) implies that either \(x \in F\) or \(y \in F\) . A prefilter with respect to the ordering opposite to that of X is called a coprefilter on X.
\end{enumerate}

\begin{enumerate}
\def\labelenumi{\alph{enumi})}
\item
  The set of upper bounds of a subset A of X which does not consist only of the least element of X is a prefilter. Give examples of prefilters which are not sets of upper bounds.
\item
  Show that if F is a prime prefilter, then CF is a prime coprefilter.
\item
  Show that if X has a least element 0, then every prefilter is contained in a maximal prefilter (use Zorn's lemma, noting that 0 belongs to no prefilter).
\item
  Find all prefilters on a linearly ordered set X, and show that they are all prime. Hence give examples of prime prefilters which are not maximal.
\item
  Let X be an ordered set of 5 elements, denoted by 0, 1, a, b, c, the order relations being 0 \textless{} a \textless{} 1, 0 \textless{} b \textless{} 1, 0 \textless{} c \textless{} 1. Show that X is a lattice and that the prefilters other than \(\{1\}\) are maximal but not prime.
\item
  Let X be a lattice having a least element 0. Show that if \(F \subset X\) is a prefilter and a is an element of X which does not belong to F, then there is a prefilter containing F and a if and only if \(\inf (a, x) \neq 0\) for all \(x \in F\) ; the smallest prefilter containing F and a is then the set of all \(y \in X\) such that \(y \geqslant \inf (a, x)\) for at least one element x of F.
\end{enumerate}

¶ 18) Let X be a distributive lattice (Set Theory, Chapter III, § 1, Exercise 16) which has a least element 0.

\begin{enumerate}
\def\labelenumi{\alph{enumi})}
\item
  If \(a\) is an irreducible element of \(\mathbf{X}\) (Set Theory, Chapter III, § 4, Exercise 7), show that the set of all upper bounds of \(a\) is a prime prefilter. Give an example of a prime prefilter which is not a set of upper bounds of an irreducible element of \(\mathbf{X}\) {[}cf.~Exercise 17 d){]}.
\item
  Show that every maximal prefilter on X is prime {[}use Exercise 17f){]}. Give examples of prime non-maximal prefilters on a distributive lattice {[}Exercise 17d){]}.
\item
  Show that every prefilter F is the intersection of the prime prefilters which contain F {[}if \(a \notin F\) , show that a maximal element U in the set of prefilters which contain F but not \(a\) is prime, by using Exercise 17f){]}.
\item
  Let \(\Omega\) be the set of prime prefilters on X. To each \(x \in \mathbf{X}\) let correspond the subset \(\mathbf{A}_x\) of \(\Omega\) consisting of prime prefilters F such that \(x \in \mathbf{F}\) . Show that the mapping \(x \to \mathbf{A}_x\) of X into \(\mathfrak{P}(\Omega)\) is injective and order-preserving, that \(\mathbf{A}_{\mathrm{inf}(x,y)} = \mathbf{A}_x \cap \mathbf{A}_y\) and that
\end{enumerate}

\[
\mathbf {A} _ {\sup (x, y)} = \mathbf {A} _ {x} \cup \mathbf {A} _ {y}.
\]

\begin{enumerate}
\def\labelenumi{\arabic{enumi})}
\setcounter{enumi}{18}
\item
  Let X be a lattice which has a least element. Show that if every prefilter on X is the intersection of the prime prefilters which contain it, then X is distributive (if the mapping \(x \rightarrow A_{x}\) is defined as in Exercise 18 d), this mapping still has the properties stated there).
\item
  A lattice X is said to be a Boolean algebra if it is distributive, has a least element 0, and a greatest element 1 and if, for each \(x \in \mathbf{X}\) , there is an element \(x' \in \mathbf{X}\) such that \(\inf (x, x') = 0\) and \(\sup (x, x') = 1\) ; such an element is said to be a complement of \(x\) in \(\mathbf{X}\) .
\end{enumerate}

\begin{enumerate}
\def\labelenumi{\alph{enumi})}
\item
  Let \(x \to \mathbf{A}_x\) be the injection of a Boolean algebra \(\mathbf{X}\) into the set of subsets of the set \(\Omega\) of prime prefilters on \(\mathbf{X}\) , defined in Exercise 18 d). Show that if \(x'\) is a complement of \(x\) in \(\mathbf{X}\) , then \(\mathbf{A}_{x'} = \complement \mathbf{A}_x\) in \(\Omega\) ; deduce that an element \(x\) has only one complement \(x'\) , that the complement of \(x'\) is \(x\) , and that the complement of \(\inf (x, y)\) is \(\sup (x', y')\) ; give direct proofs of these statements. The ideals of the Boolean ring \(\mathbf{A}\) corresponding to \(\mathbf{X}\) , other than \(\mathbf{A}\) , are the coprefilters on \(\mathbf{X}\) .
\item
  Show that, in a Boolean algebra \(\mathbf{X}\) , every prime prefilter is maximal (remark that, for each \(x \in \mathbf{X}\) , every prime prefilter must contain either \(x\) or its complement).
\item
  Conversely, let X be a distributive lattice which has a least element 0 and a greatest element 1, and is such that every prime prefilter on X is maximal. Show that X is a Boolean algebra. (For an element \(x \in X\) , consider the coprefilter C (Exercise 17) consisting of those \(y \in X\) such that \(\inf(x, y) = 0\) ; if \(x\) has no complement, there is a maximal coprefilter M containing \(x\) and C; remark that the complement U of M in X is a prime prefilter {[}Exercise 17 b) and 18 b){]} and use Exercise 17 f)). d) Show that if X is a topological space then the set B of subsets of X which are both open and closed is a Boolean algebra if ordered by inclusion. If we take X to be the topological space associated with the Fréchet filter (no. 5), show that the Boolean algebra B thus obtained is countably infinite and therefore cannot be isomorphic to the ordered set P(A) of all subsets of a set A.
\end{enumerate}
% semantic-fragments: legacy-input-seam
\relax{}
